\section{A correction with small Fourier coefficients}
\label{sec:correction}

The function furnished by Proposition~\ref{prop:starting} can be small on
short intervals.  We now add a wave of fixed amplitude on each such interval.
Its phase has two parts: $N$ times a fixed piecewise quadratic leading
phase, and an affine adjustment whose slope is bounded independently
of $N$. Where the correction meets a nonzero original wave, we match
their leading phase derivatives and use the adjustment to match their
full phase values modulo integers at the endpoint. Turning on the amplitude in a short
strip then prevents cancellation.

On most of each gap, the leading phase derivative runs through a range
reserved for that gap and its reflection. On these traversals, a given
Fourier coefficient can receive stationary contributions from at most
one gap and its reflection. In the oscillatory integrals, we keep the affine adjustment
in the amplitude; it is the leading phase derivative that governs the
separation.

\begin{proposition}[Correction on short gaps]\label{prop:correction}
Fix $0<\delta<1/100$, and put
\[
 b=\frac18,
 \qquad K_\delta=\sqrt{\frac{(1+\delta)^3}{1-\delta}}.
\]
Suppose that $B_N\in C^\infty(\T;\C)$, $N\in\N$, satisfy
$B_N(-t)=\overline{B_N(t)}$, and that their common modulus
$B=|B_N|\in C^\infty(\T)$ has the following properties.
\begin{enumerate}
\item $0\le B\le K_\delta$, and $B$ vanishes in neighborhoods of $0$ and
  $1/2$.  The set
  \[
    E=\{t\in[0,1/2]:B(t)\ge b\}
  \]
  is a nonempty finite union of disjoint closed intervals of positive length.
  The closures of its complementary intervals are
  $J_j=[u_j,v_j]$, $1\le j\le M$, with lengths $l_j=v_j-u_j>0$, and
  \[
    0<L:=\sum_{j=1}^M l_j\le 10\delta.
  \]
\item At every endpoint $z\in(0,1/2)$ of a gap $J_j$, there are a
  neighborhood of $z$, a real number $\alpha_z$, and a real quadratic
  polynomial $\psi_z$, all independent of $N$, such that
  \[
    B_N(t)=B(t)\e\bigl(\alpha_z+N\psi_z(t)\bigr),
    \qquad x_z:=\psi_z'(z)\in(0,1).
  \]
\end{enumerate}
There is an absolute constant $D$ such that, for every sufficiently large
integer $N$, there is a continuous function $R_N:\T\to\C$ satisfying
$R_N(-t)=\overline{R_N(t)}$, supported on the gaps and their reflections,
for which
\begin{align}
 b-o(1)&\le |B_N(t)+R_N(t)|\le K_\delta
       &&(t\in\T),\label{eq:corrected-envelope}\\
 \max_{0\le k<N}\sqrt N\,|\widehat R_N(k)|
       &\le D\sqrt\delta,\label{eq:correction-coefficients}\\
 \sum_{k\in\Z\setminus\{0,\ldots,N-1\}}|\widehat R_N(k)|
       &=O_\delta(N^{-1/4}).\label{eq:correction-tail}
\end{align}
Every Fourier coefficient $\widehat R_N(k)$ is real.  The threshold for $N$
and the constants implicit in $o(1)$ and $O_\delta$ may depend on all the
fixed input data; $D$ does not.
\end{proposition}

The hypotheses are precisely the geometric and local phase properties
established in Proposition~\ref{prop:starting}.  We prove the proposition
by first constructing the correction and checking its modulus, and then
estimating its Fourier coefficients.  All choices made below are fixed
before $N$ tends to infinity, except the explicitly indexed amplitudes and
phase adjustments.

\subsection{Assigning the phase derivatives}

Call the gap endpoints in $(0,1/2)$ \emph{transition endpoints}.  At these
points $B(z)=b$, and the prescribed derivative is $x_z=\psi_z'(z)$.
At the two remaining endpoints set $x_0=x_{1/2}=1/8$.
For each gap choose a closed interval
\[
 D_j=[p_j,q_j]\subset(1/4,3/4),
 \qquad s_j:=q_j-p_j=\frac{l_j}{4L},
\]
so that these intervals are pairwise disjoint and none of their endpoints
equals any prescribed $x_z$.  We call $D_j$ the derivative slot of $J_j$.
Such a choice is possible because $\sum_j s_j=1/4$: first place the
intervals with positive separations inside $(1/4,3/4)$, and then translate
them together by a sufficiently small amount avoiding the finitely many
forbidden endpoint equalities.

We will make the leading derivative traverse $D_j$ linearly on a subinterval of
length at most $l_j$, giving curvature at least $s_j/l_j=1/(4L)$.
For amplitudes with uniformly bounded supremum norm and total variation,
the quadratic estimate proved below
then gives a coefficient cost $O(\sqrt{L/N})$ for this traversal and its
reflection. The separated slots allow at most one such pair to be
stationary for a given frequency. Joining these traversals to the
prescribed endpoint derivatives may cross other ranges; we will control
those contributions separately by using short intervals with large curvature.

Choose positive lengths $h_j^-,h_j^+<l_j/4$.  Define a continuous,
piecewise affine function $\xi_j:J_j\to(0,1)$ by the three successive
linear interpolations
\[
\begin{array}{c|c}
\text{interval}&\text{endpoint values of }\xi_j\\ \hline
[u_j,u_j+h_j^-]&x_{u_j},\ p_j\\[2pt]
[u_j+h_j^-,v_j-h_j^+]&p_j,\ q_j\\[2pt]
[v_j-h_j^+,v_j]&q_j,\ x_{v_j}.
\end{array}
\]
We refer to the first and last intervals as the outer pieces and to the
remaining interval as the middle piece.  The slopes on the outer pieces
are nonzero.  As their lengths tend to zero, their absolute slopes tend
to infinity.  We may therefore choose these lengths so that
\begin{equation}\label{eq:outer-slope-budget}
 \sum_{\substack{\text{outer pieces}\\\text{in }[0,1/2]}}
       |\xi_j'|^{-1/2}\le\sqrt\delta.
\end{equation}
There are only finitely many terms in this sum.  The slope on every middle
piece satisfies
\begin{equation}\label{eq:middle-slope}
 d_j:=\frac{s_j}{l_j-h_j^--h_j^+}\ge\frac1{4L}.
\end{equation}
Let $\phi_j$ be a primitive of $\xi_j$.  Thus $\phi_j$ is $C^1$ on $J_j$
and quadratic on each of its three pieces.  Since all the endpoint values
used in the interpolations lie in $(0,1)$, there is a fixed
$c_\delta>0$ such that
\begin{equation}\label{eq:derivative-interior}
 c_\delta\le\phi_j'(t)\le1-c_\delta
       \qquad(t\in J_j,\ 1\le j\le M).
\end{equation}
Here and below, a subscript $\delta$ allows dependence on all the fixed
data.  This compact containment will keep the Fourier mass of the
correction inside the coefficient interval $0\le k<N$.

\subsection{Matching phases while turning on the correction}

Set $A=1/2$.  For each $N$, choose the endpoint values of an affine function
$\gamma_{j,N}:J_j\to\R$ in $[0,1)$ so that
\begin{equation}\label{eq:endpoint-phase-match}
 \e\bigl(N\phi_j(z)+\gamma_{j,N}(z)\bigr)=
 \begin{cases}
   B_N(z)/b,&z\text{ is a transition endpoint},\\
   1,&z\in\{0,1/2\}.
 \end{cases}
\end{equation}
Any point of the unit circle has a representative with argument in
$[0,1)$, so both endpoint conditions can be imposed independently.  Their
affine interpolation satisfies
\begin{equation}\label{eq:gamma-bounds}
 |\gamma_{j,N}'|\le l_j^{-1},
 \qquad \int_{J_j}|\gamma_{j,N}'(t)|\,dt\le1.
\end{equation}

We turn the amplitude on over strips of width $w=w_N$.  The requirements
$Nw^2\to0$ and $(Nw)^{-1}\to0$ will serve different purposes: the first
preserves phase alignment on a strip, and the second makes the exterior
Fourier tail tend to zero.  We take
\begin{equation}\label{eq:taper-width}
 w=N^{-3/4}.
\end{equation}
Fix a smooth nondecreasing function $\vartheta:\R\to[0,1]$ with $\vartheta=0$ on
$(-\infty,0]$ and $\vartheta=1$ on $[1,\infty)$.  On $J_j$ define
$r_{j,N}$ as the product of the following factors: include
$\vartheta((t-u_j)/w)$ if $u_j$ is a transition endpoint, and include
$\vartheta((v_j-t)/w)$ if $v_j$ is a transition endpoint.  No factor is included
at $0$ or $1/2$.  For large enough $N$ the two strips are disjoint, each
lies within its outer piece, and each lies in the neighborhood where the
corresponding local formula for $B_N$ holds.  Thus $r_{j,N}=1$ off the
transition strips, it vanishes to all orders at transition endpoints, and
\begin{equation}\label{eq:taper-variation}
 0\le r_{j,N}\le1,
 \qquad \int_{J_j}|r_{j,N}'(t)|\,dt\le2.
\end{equation}

On the positive half-circle put
\begin{equation}\label{eq:correction-definition}
 R_N(t)=
 \begin{cases}
 A r_{j,N}(t)\e\bigl(N\phi_j(t)+\gamma_{j,N}(t)\bigr),&t\in J_j,\\
 0,&t\notin\bigcup_j J_j,
 \end{cases}
\end{equation}
and extend it by $R_N(-t)=\overline{R_N(t)}$.  At transition endpoints the
amplitude is zero, while \eqref{eq:endpoint-phase-match} gives
\begin{equation}\label{eq:reflection-values}
 R_N(0)=R_N(1/2)=A.
\end{equation}
The extension is therefore continuous both at $0$ and at the identified
points $-1/2$ and $1/2$.  In particular, the two points corresponding to
$z=1$ and $z=-1$ on the unit circle are covered by the construction.
Figure~\ref{fig:gap-slots} separates the fixed outer pieces from the
shrinking taper strips and illustrates the use of disjoint derivative
slots.

\begin{figure}[tbp]
\centering
\begin{tikzpicture}[
  x=1cm,y=1cm,
  font=\small,
  >=Stealth,
  phase/.style={draw=blue!65!black,line width=1.15pt},
  taper/.style={draw=orange!60!black,line width=1.1pt},
  guide/.style={draw=black!35,densely dashed,line width=.4pt},
  axis/.style={draw=black!75,line width=.5pt,->},
  slot/.style={draw=blue!65!black,line width=3pt},
  neighborhood/.style={draw=black!55,densely dashed,fill=black!3,line width=.5pt}
]
\node[anchor=west,font=\small\bfseries] at (0,6.55)
  {(a) One interior gap $J_j=[u_j,v_j]$};
\fill[blue!3] (1.4,4.1) rectangle (8.4,5.0);
\fill[orange!13] (1.4,.65) rectangle (1.65,5.9);
\fill[orange!13] (8.15,.65) rectangle (8.4,5.9);
\draw[axis] (.9,3.2) -- (.9,5.95);
\draw[axis] (1.15,3.2) -- (8.95,3.2) node[right] {$t$};
\node[anchor=west] at (.65,6.15) {$\xi_j=\phi_j'$};
\draw[guide] (.9,4.1) -- (2.2,4.1);
\draw[guide] (.9,5.0) -- (7.6,5.0);
\node[anchor=east] at (.82,4.1) {$p_j$};
\node[anchor=east] at (.82,5.0) {$q_j$};
\draw[guide] (2.2,3.2) -- (2.2,4.1);
\draw[guide] (7.6,3.2) -- (7.6,5.0);
\draw[black!85,line width=1.15pt] (1.4,5.65) -- (2.2,4.1);
\draw[phase] (2.2,4.1) -- (7.6,5.0);
\draw[black!85,line width=1.15pt] (7.6,5.0) -- (8.4,3.9);
\foreach \x/\y in {1.4/5.65,2.2/4.1,7.6/5.0,8.4/3.9}
  \fill[black!85] (\x,\y) circle (1.35pt);
\node[anchor=west,inner sep=2pt] at (1.55,5.6) {$x_{u_j}$};
\node[anchor=west,inner sep=2pt] at (8.4,3.9) {$x_{v_j}$};
\node[anchor=south,text=blue!65!black] at (4.95,5.18)
  {$\xi_j'=d_j\ge 1/(4L)$};
\draw (1.4,3.26) -- (1.4,3.14);
\draw (8.4,3.26) -- (8.4,3.14);
\node[anchor=north] at (1.4,3.13) {$u_j$};
\node[anchor=north] at (8.4,3.13) {$v_j$};
\draw[<->,black!65] (1.4,2.6) -- (2.2,2.6)
  node[midway,below=2pt] {$h_j^-$};
\draw[<->,black!65] (2.2,2.6) -- (7.6,2.6)
  node[midway,below=2pt] {middle piece};
\draw[<->,black!65] (7.6,2.6) -- (8.4,2.6)
  node[midway,below=2pt] {$h_j^+$};
\draw[axis] (.9,.65) -- (.9,1.95);
\draw[axis] (1.15,.65) -- (8.95,.65) node[right] {$t$};
\node[anchor=east] at (.82,1.73) {$1$};
\node[anchor=east] at (.82,.65) {$0$};
\node[text=orange!60!black] at (4.9,1.15) {$r_{j,N}$};
\draw[taper] (1.4,.65)
  .. controls (1.535,.65) and (1.515,1.73) .. (1.65,1.73)
  -- (8.15,1.73)
  .. controls (8.285,1.73) and (8.265,.65) .. (8.4,.65);
\draw[guide] (1.65,.65) -- (1.65,1.73);
\draw[guide] (8.15,.65) -- (8.15,1.73);
\draw[orange!60!black] (1.4,.30) -- (1.65,.30);
\draw[orange!60!black] (8.15,.30) -- (8.4,.30);
\foreach \x in {1.4,1.65,8.15,8.4}
  \draw[orange!60!black] (\x,.24) -- (\x,.36);
\node[anchor=north] at (1.525,.20) {$w$};
\node[anchor=north] at (8.275,.20) {$w$};
\node[anchor=west,font=\small\bfseries] at (10.05,6.55)
  {(b) Derivative slots};
\draw[axis] (10.4,1.6) -- (10.4,6.0);
\node[anchor=west] at (10.75,6.05) {derivative value};
\foreach \y/\lab in {1.6/0,2.675/{\frac14},4.825/{\frac34},5.9/1} {
  \draw (10.34,\y) -- (10.46,\y);
  \node[anchor=east] at (10.28,\y) {$\lab$};
}
\draw[neighborhood] (11.15,2.71) rectangle (11.75,3.26);
\draw[neighborhood] (11.15,3.36) rectangle (11.75,4.16);
\draw[neighborhood] (11.15,4.26) rectangle (11.75,4.825);
\draw[slot] (11.45,2.85) -- (11.45,3.12);
\draw[slot] (11.45,3.5) -- (11.45,4.02);
\draw[slot] (11.45,4.4) -- (11.45,4.685);
\node[anchor=west] at (11.98,2.985) {$D_a$};
\node[anchor=west] at (11.98,3.76) {$D_j=[p_j,q_j]$};
\node[anchor=west] at (11.98,4.5425) {$D_b$};
\draw[black!75,densely dotted] (10.4,4.11) -- (14.65,4.11);
\fill (11.45,4.11) circle (1.5pt);
\node[anchor=west] at (14.7,4.11) {$k/N$};
\node[align=center,text width=4.7cm] at (12.65,1.0)
  {At most one slot is close to $k/N$.};
\end{tikzpicture}
\caption{Schematic of the correction.  In (a), the leading phase
derivative $\xi_j=\phi_j'$ matches the prescribed endpoint values and
traverses $D_j$ on the middle piece.  The outer lengths $h_j^\pm$ are
fixed; the shaded taper strips have shrinking width $w=N^{-3/4}$,
enlarged here.  The lower curve shows the amplitude factor for an
interior gap.  In (b), the dashed neighborhoods of the slots are fixed
and disjoint.  At most one middle piece and its reflection can have a
leading phase derivative close to $k/N$; all other middle pieces are
nonstationary.}
\label{fig:gap-slots}
\end{figure}


On a gap outside its transition strips, the reverse triangle inequality
gives
\[
 |B_N+R_N|\ge A-b=\frac38.
\]
To check a transition strip at $z$, compare the phases after subtracting
their values at $z$.  Since $\phi_j'(z)=\psi_z'(z)$ and both derivatives
are Lipschitz near $z$,
\[
 \bigl|\phi_j(t)-\phi_j(z)-\psi_z(t)+\psi_z(z)\bigr|
       \le C_\delta|t-z|^2.
\]
Together with \eqref{eq:gamma-bounds} and the endpoint phase agreement,
this shows that the two phases differ, modulo an integer, by at most
\begin{equation}\label{eq:phase-strip-error}
 C_\delta(Nw^2+w)=O_\delta(N^{-1/2}).
\end{equation}
Consequently, throughout that strip,
\[
 B_N(t)+R_N(t)
 =\e\bigl(\alpha_z+N\psi_z(t)\bigr)
       \bigl(B(t)+A r_{j,N}(t)+O_\delta(N^{-1/2})\bigr).
\]
Because $B(z)=b$ and $B$ is fixed and smooth, $B(t)\ge b-O_\delta(w)$
there.  This proves the lower bound $b-o(1)$ on every strip.  Off the gaps,
$R_N=0$ and $B\ge b$, and reflection preserves all these modulus bounds.
The upper bound is simpler: on gaps $|B_N+R_N|\le b+A=5/8<K_\delta$,
and elsewhere it is at most $K_\delta$.  This proves
\eqref{eq:corrected-envelope}.

The correction now gives the required lower envelope.  The remaining
task is to show that its Fourier coefficients cost only $O(\sqrt\delta)$
on the scale $N^{-1/2}$ and that discarding frequencies outside
$\{0,\ldots,N-1\}$ changes it uniformly by $o(1)$.

\subsection{A bound for each Fourier coefficient}

Split the integral for $\widehat R_N(k)$ into the three pieces of every
gap and their reflections.  On a positive piece its contribution has
the form
\begin{equation}\label{eq:piece-integral}
 \int_I a(t)\e\bigl(N\phi(t)-kt\bigr)\,dt,
 \qquad a(t)=A r_{j,N}(t)\e\bigl(\gamma_{j,N}(t)\bigr),
 \quad \phi=\phi_j|_I.
\end{equation}
On the reflected piece we use
$\phi(t)=-\phi_j(-t)$ and
$a(t)=\overline{A r_{j,N}(-t)\e(\gamma_{j,N}(-t))}$.
In particular, the reflected phase derivative is
\begin{equation}\label{eq:reflected-derivative}
 \phi'(t)=\xi_j(-t).
\end{equation}
Thus the original and reflected middle pieces use the same slot $D_j$.
By \eqref{eq:gamma-bounds} and \eqref{eq:taper-variation}, on every piece
\begin{equation}\label{eq:amplitude-bv}
 \sup_I|a|+\int_I|a'(t)|\,dt
 \le A+2A+2\pi A=:C_1.
\end{equation}
This constant is absolute, even if a gap is very short or there are many
gaps.

We use two elementary oscillatory integral estimates.  If $\phi$ is
quadratic with $\phi''=d\ne0$, then
\begin{equation}\label{eq:quadratic-integral}
 \left|\int_I a(t)\e\bigl(N\phi(t)-kt\bigr)\,dt\right|
 \le \frac{C}{\sqrt{N|d|}}
       \left(\sup_I|a|+\int_I|a'|\right).
\end{equation}
Indeed, completing the square and scaling reduce the integral without
$a$ over any subinterval to an integral of $\e(\pm u^2/2)$, multiplied
by $(N|d|)^{-1/2}$.  Those integrals are uniformly bounded: on $[-1,1]$
use length, and on each exterior half-line integrate by parts using the
phase derivative $\pm u$.  Integration against a primitive then proves
\eqref{eq:quadratic-integral}.  We also use the first-derivative bound
\begin{equation}\label{eq:first-derivative-integral}
 \left|\int_I a(t)\e\bigl(N\phi(t)-kt\bigr)\,dt\right|
 \le \frac{C}{N\rho}
       \left(\sup_I|a|+\int_I|a'|\right)
\end{equation}
when $N\phi'-k$ is monotone and has absolute value at least $N\rho$ on
$I$.  One integration by parts proves this bound, since the total
variation of $1/(N\phi'-k)$ is at most $2/(N\rho)$.

Let $C_2$ be the absolute constant obtained by combining
\eqref{eq:amplitude-bv} with \eqref{eq:quadratic-integral}.
The outer pieces and their reflections contribute at most
\[
 \frac{2C_2\sqrt\delta}{\sqrt N}
\]
by \eqref{eq:outer-slope-budget}.  For the middle pieces choose a fixed
$\rho>0$ so that the closed $\rho$-neighborhoods of the finitely many
slots $D_j$ are disjoint.  For any $k$, at most one slot is within
distance $\rho$ of $k/N$.  On that middle piece and its reflection,
\eqref{eq:quadratic-integral} and \eqref{eq:middle-slope} give a total
bound
\[
 \frac{4C_2\sqrt L}{\sqrt N}.
\]
On all remaining middle pieces, $|N\phi'-k|\ge N\rho$, so
\eqref{eq:first-derivative-integral} bounds their total contribution
by $O_\delta(N^{-1})$, uniformly in $k$.  The number of pieces and the
separation $\rho$ enter only this last error.  We conclude that
\[
 \sqrt N\,|\widehat R_N(k)|
 \le 2C_2\sqrt\delta+4C_2\sqrt L+O_\delta(N^{-1/2}).
\]
Since $L\le10\delta$, for sufficiently large $N$ this is at most
$D\sqrt\delta$ with, for example,
$D=2C_2+4\sqrt{10}C_2+1$.  This proves
\eqref{eq:correction-coefficients} with an absolute constant.
Hermitian symmetry also gives
\[
 \overline{\widehat R_N(k)}
 =\int_\T\overline{R_N(t)}\e(kt)\,dt
 =\int_\T R_N(-t)\e(kt)\,dt
 =\widehat R_N(k).
\]

\subsection{The Fourier mass outside the coefficient interval}

Let $k<0$ or $k\ge N$, and write $T=N+|k|$.
On each piece in \eqref{eq:piece-integral}, set
\[
 V(t)=N\phi'(t)-k.
\]
The derivative containment \eqref{eq:derivative-interior}, also valid on
reflected pieces by \eqref{eq:reflected-derivative}, implies
\begin{equation}\label{eq:exterior-derivatives}
 |V(t)|\ge c_\delta T,
 \qquad |V'(t)|\le C_\delta N,
 \qquad V''(t)=0
\end{equation}
on each piece interior, after decreasing $c_\delta$ if necessary.
The scaled smooth steps and the affine phase adjustments give
\begin{equation}\label{eq:amplitude-two-derivatives}
 \sup_I|a'|\le C_\delta w^{-1},
 \qquad \int_I|a''(t)|\,dt\le C_\delta w^{-1}.
\end{equation}
For the second estimate, the term of size $w^{-2}$ from a differentiated
step is supported on a strip of length $w$; the other terms are bounded
using \eqref{eq:gamma-bounds}.  Recall also the absolute bounds on
$\sup|a|$ and $\int|a'|$ in \eqref{eq:amplitude-bv}.

One integration by parts on a piece gives
\begin{equation}\label{eq:first-ibp-tail}
 \int_I a\,\e(N\phi-kt)
 =\left[\frac{a\,\e(N\phi-kt)}{2\pi iV}\right]_{\partial I}
 -\frac1{2\pi i}\int_I(a/V)'\,\e(N\phi-kt).
\end{equation}
The boundary terms in this first integration cancel when all pieces
are summed.  At a join within a gap, the numerator is the continuous
function $R_N(t)\e(-kt)$ and the denominator agrees because $\phi_j'$
is continuous.  At transition endpoints the numerator vanishes.
At $0$, the reflected and original phase derivatives both equal $x_0$,
and the numerator has the common value $A$.  Finally, at the identified
endpoints $-1/2$ and $1/2$, both phase derivatives equal $x_{1/2}$ and
the numerator has the common value $A(-1)^k$.  Thus cancellation holds
on the full circle, including the two reflection joins.

For a second integration by parts define, on each piece,
\[
 Q=\frac{(a/V)'}{V}
   =\frac{a'}{V^2}-\frac{aV'}{V^3}.
\]
No cancellation of the new boundary terms is needed.  Using $V''=0$,
\begin{equation}\label{eq:second-ibp-algebra}
 Q'=\frac{a''}{V^2}
       -\frac{3a'V'}{V^3}
       +\frac{3a(V')^2}{V^4}.
\end{equation}
Equations \eqref{eq:exterior-derivatives} and
\eqref{eq:amplitude-two-derivatives} bound the endpoint values of $Q$ by
\[
 C_\delta\left(\frac{w^{-1}}{T^2}+\frac{N}{T^3}\right)
\]
and the integral of $|Q'|$ by
\[
 C_\delta\left(\frac{w^{-1}}{T^2}
                 +\frac{N}{T^3}+\frac{N^2}{T^4}\right).
\]
Summing over the fixed finite number of pieces in
\eqref{eq:first-ibp-tail}, and integrating its remaining integrals once
more, therefore gives
\[
 |\widehat R_N(k)|
 \le C_\delta\left(\frac{w^{-1}}{T^2}
                 +\frac{N}{T^3}+\frac{N^2}{T^4}\right)
 \le \frac{C_\delta w^{-1}}{(N+|k|)^2}.
\]
Since
$\sum_{k<0\text{ or }k\ge N}(N+|k|)^{-2}=O(N^{-1})$,
we obtain
\[
 \sum_{k<0\text{ or }k\ge N}|\widehat R_N(k)|
       \le C_\delta(Nw)^{-1}=O_\delta(N^{-1/4}),
\]
as asserted in \eqref{eq:correction-tail}.  This completes the proof
of Proposition~\ref{prop:correction}.
\hfill$\square$

In particular, the Fourier series of $R_N$ is absolutely convergent.
Its sum equals $R_N$: both are continuous and have the same Fourier
coefficients, so uniqueness of Fourier coefficients applies.  Hence
\begin{equation}\label{eq:correction-projection}
 \sup_{t\in\T}
 \left|R_N(t)-\sum_{k=0}^{N-1}\widehat R_N(k)\e(kt)\right|
       =O_\delta(N^{-1/4}).
\end{equation}
We may therefore add the correction using only the required consecutive
degrees, with a uniform error tending to zero.
